
\part {기초}

\chapter{\simulator 을 이용한 건물 에너지 시뮬레이션의 목적}

쓸 말 없으면 삭제

\chapter{건물 모델링 workflow}

\section{일반사항}

\begin{itemize}
    \item 모든 부위별 성능 값은 승인된 설계도서(도면, 시방서 및 시험성적서)를 기준으로 입력하는 것을 원칙으로 한다. 설계도서에서 성능 값을 확인할 수 없는 경우 본 문서의 기준에 따른다.
    \item 설계도서를 확인할 수 없으면서 입력하지 않아도 되는 변수는 입력하지 않고 시뮬레이션 엔진의 기본값을 사용한다.
    \item 본 문서에 기재되지 않은 내용은 ``건축물의 에너지효율등급 평가''에서의 판단 기준에 준하여 적용한다. 
\end{itemize}

\section{입력 데이터의 구조, 건물의 분석 및 입력 순서}
다이어그램들..

권장되는 순서..